\documentclass[11pt]{article}
\usepackage[T1]{fontenc}
\usepackage[margin=0.86in]{geometry}
\usepackage{amsmath,amssymb,amsthm,booktabs,array,hyperref}
\hypersetup{hidelinks}
\newtheorem{theorem}{Theorem}
\newcommand{\Q}{\mathbb Q}
\newcommand{\Gal}{\operatorname{Gal}}
\newcommand{\Frob}{\operatorname{Frob}}
\newcommand{\Tr}{\operatorname{tr}}
\newcommand{\sgn}{\operatorname{sgn}}
\title{A split-prime test for the cubic Kummer \(L\)-function at 269}
\author{AI-created mathematical note}
\date{22 November 2026}
\begin{document}
\maketitle
\noindent
There is a single prime-counting error term which would settle the cubic
example left open by a rational Dirichlet trace formula. The example also
admits an exact calculation of how much local Frobenius information is
visible to rational characters. Both calculations follow from classical
Artin formalism and Chebotarev's theorem. The prime-counting estimate itself
remains an additional hypothesis.

\section{The cubic test}
Write \(F=\Q(\sqrt[3]{269})\), \(K=\Q(\omega)\), where
\(\omega^2+\omega+1=0\), and \(E=FK\). For unramified rational primes
\(p\), set
\[
 \theta_{\rm sp}(x)=
 \sum_{\substack{p\le x,\ p\nmid 3\cdot269\\
 p\equiv1\bmod3,\ 269^{(p-1)/3}\equiv1\bmod p}}\log p .
\]
The congruence inside the sum tests complete splitting in \(E\).
Indeed \(X^3-269\) is separable at these primes, \(p\equiv1\bmod3\)
puts all cube roots of unity in the residue field, and the displayed
power test is precisely the cube criterion there.

\begin{theorem}\label{thm:test}
Assume RH for \(\zeta_{\Q}\) and for the primitive quadratic
\(L(s,\chi_{-3})\). Let \(\eta\) be either nontrivial cubic character of
\(\Gal(E/K)\), viewed as a finite-order Hecke character of \(K\). These
three statements are equivalent:
\begin{enumerate}
\item All zeros of \(L_K(s,\eta)\) in \(0<\Re s<1\) satisfy \(\Re s=1/2\).
\item \(\theta_{\rm sp}(x)=x/6+O(\sqrt{x}\log^2(2x))\).
\item For every \(\epsilon>0\),
\(\theta_{\rm sp}(x)=x/6+O_\epsilon(x^{1/2+\epsilon})\).
\end{enumerate}
\end{theorem}
\begin{proof}
The Eisenstein cubic has nonsquare discriminant \(-27\cdot269^2\), so its
splitting group is \(S_3\) and its quadratic resolvent is \(K\). Inducing a
nontrivial character of \(A_3\) gives the standard irreducible representation
\(V\) of \(S_3\). The two conjugate cubic Hecke functions coincide under
the conjugation of ideals of \(K\). The regular character and Artin formalism give
\[
 \zeta_E(s)=\zeta_K(s)L_K(s,\eta)^2,\qquad
 \zeta_K(s)=\zeta(s)L(s,\chi_{-3}).
\]
The Hecke function is entire and all factors are holomorphic in the open
strip. Hence the two assumed RH statements make (1) equivalent to RH for
\(\zeta_E\). Under this last RH, the prime-ideal explicit formula
(for example Lagarias--Odlyzko, Theorem~9.1 applied over \(E\))
gives \(\psi_E(x)=x+O_E(\sqrt{x}\log^2(2x))\), where \(\psi_E\)
counts prime-ideal powers with weight \(\log N\mathfrak p\).

For this fixed Galois field,
\[
 \psi_E(x)=6\theta_{\rm sp}(x)+O_E(\sqrt{x}\log^2(2x)).
\]
In fact a rational unramified prime has a degree-one place precisely when
it splits completely, in which case it has six. For contributions of norm
\(p^j\) with \(j\ge2\), each \((p,j)\) has total weight at most
\(6\log p\). Summing the elementary bounds
\(\sum_{p\le x^{1/j}}\log p\le x^{1/j}\log x\) up to \(j\le\log_2x\)
suffices. The two excluded rational primes cost \(O_E(\log x)\).
This proves (1) implies (2), which implies (3).

Under (3), \(R(x)=\psi_E(x)-x=O_{\epsilon,E}(x^{1/2+\epsilon})\).
For \(\Re s>1\), partial summation yields
\[
 -\frac{\zeta'_E}{\zeta_E}(s)=\frac{s}{s-1}
                  +s\int_1^\infty R(x)x^{-s-1}\,dx .
\]
The integral is holomorphic throughout \(\Re s>1/2\). Thus \(\zeta_E\)
has no zero there. Its functional equation gives the other half of RH.
\end{proof}

\section{What rational residue classes see}
At a good prime put
\[
 a(p)=\#\{x\bmod p:x^3=269\}-1=\Tr V(\Frob_p).
\]
Table~\ref{tab:classes} describes all three possible Frobenius classes.
For example \(13\) and \(37\) are both \(1\bmod12\). The first has no
root and the second has roots \(7,33,34\).

\begin{table}[ht]
\centering
\begin{tabular}{llrrr}
\toprule
Frobenius & Factor degrees & Prime density & \(a(p)\) & \(\sgn(\Frob_p)\)\\
\midrule
Identity & \(1,1,1\) & \(1/6\)& \(2\)&\(1\)\\
Transposition & \(1,2\)& \(1/2\)& \(0\)&\(-1\)\\
Three-cycle & \(3\)& \(1/3\)& \(-1\)&\(1\)\\
\bottomrule
\end{tabular}
\caption{The action on the three cubic roots gives the factor degrees and
the standard character. The sign is \(\chi_{-3}(p)\).}\label{tab:classes}
\end{table}

\begin{theorem}\label{thm:residues}
For every fixed modulus \(m\) and every complex function \(b\) on its
reduced residue classes,
\[
 \lim_{x\to\infty}\frac1{\pi(x)}
   \sum_{\substack{p\le x\\p\nmid 3\cdot269m}}
   |a(p)-b(p\bmod m)|^2
 =1+\frac1{\varphi(m)}\sum_{u\in(\mathbb Z/m)^\times}|b(u)|^2.
\]
Across all finite moduli, the least possible density of mistakes when
predicting the exact value of \(a(p)\) from \(p\bmod m\) is \(1/6\).
\end{theorem}
\begin{proof}
Put \(B=\Q(\zeta_m)\). The Galois group of \(EB\) is the fibre product
over \(E\cap B\). This intersection is abelian over \(\Q\), hence is
\(\Q\) or \(K\). Conditional on \(\Frob_p|_B=u\), Chebotarev assigns
uniform element weights on \(S_3\) in the first case. In the second it
assigns uniform element weights on the prescribed parity coset. Each coset
has mean trace zero: the even coset contributes \(2-1-1=0\) and all
three odd values are zero. The mean square on the full group is
\((4+2)/6=1\). Expanding the square and using the uniform distribution
on \(\Gal(B/\Q)\) proves the identity.

With the parity available, every odd prime class is classified correctly
by prediction zero. In an even class two of the three Frobenius elements
have trace \(-1\), leaving conditional error \(1/3\). Half the primes
lie there. If \(E\cap B=\Q\), no residue class changes the original
probabilities and the smallest possible error is \(1/2\). Taking
\(m=3\) achieves the global minimum \(1/6\).
\end{proof}

\paragraph{Norm twists.}
The same mean-square lower bound \(1\) survives each finite linear
combination \(b(p)=\sum_j b_j(p\bmod m)p^{it_j}\), with fixed real \(t_j\).
For each \(u\) the centered sum \(A_u(x)=\sum_{p\le x,\ p\equiv u(m)}a(p)\)
is \(o(x/\log x)\) by the preceding proof and the prime number theorem.
Partial summation gives
\[
 \sum_{p\le x,\ p\equiv u(m)} a(p)p^{it}
 = x^{it}A_u(x)-it\int_2^x A_u(y)y^{it-1}dy+O(1)
 =o(x/\log x).
\]
All cross terms against the target vanish in the normalized mean. The
liminf of the squared error is at least \(1\), attained by \(b=0\).
The conclusion concerns fixed predictors as \(x\to\infty\).

\paragraph{Prime powers.}
For \(r\ge1\), the \(p^r\) term of the logarithmic derivative uses
\(\Tr V(\Frob_p^r)\). Averaging the three values in each sign coset
gives the following exact periodic table. The residual is obtained by
subtracting that average.
\[
\begin{array}{c|rrrrrr}
r\bmod6&1&2&3&4&5&0\\\hline
\text{even mean}&0&0&2&0&0&2\\
\text{odd mean}&0&2&0&2&0&2\\
\text{residual}&a(p)&a(p)&0&a(p)&a(p)&0
\end{array}
\]
To verify each column, raise the identity, a transposition, and a
three-cycle to the indicated power and read Table~\ref{tab:classes}.
In particular the third-power term is entirely rational-quadratic, while
the first-power term loses its entire mean-square mass under abelian
projection. At first powers the exact contrast
\(a(p)=3{\bf1}_{p\text{ splits in }E}-{\bf1}_{p\text{ splits in }K}\)
isolates the missing signal. Theorem~\ref{thm:test} asks for its
cancellation as the prime bound grows. Orthogonality alone supplies no
square-root error.

\paragraph{Higher finite-image ranks.}
For any nonlinear irreducible complex representation \(W\) of a finite
Galois group \(G\) over \(\Q\), its Frobenius character has the same
mean-square distance \(1\) from the best fixed abelian residue-class
predictor. Indeed \(W^{[G,G]}\) is a \(G\)-stable subspace. If nonzero,
it would be all of \(W\), making \(W\) an irreducible representation of
an abelian group and hence one-dimensional. Averaging over any coset of
a normal subgroup containing \([G,G]\) therefore has zero character
sum. Character orthogonality gives \(\langle\chi_W,\chi_W\rangle=1\).
Chebotarev in the compositum with \(\Q(\zeta_m)\) now gives the argument
of Theorem~\ref{thm:residues}. This is a statement about Frobenius data.
The zero locations of other \(L\)-functions require their own analytic input.

\paragraph{References.}
J.~C. Lagarias and A.~M. Odlyzko, \emph{Effective versions of the Chebotarev
density theorem}, in \emph{Algebraic Number Fields: \(L\)-Functions and Galois
Properties}, Academic Press, 1977, pp.~409--464, Theorem~9.1.
A.~Connes, C.~Consani and M.~Marcolli, \emph{The Weil proof and the geometry
of the adeles class space}, \href{https://arxiv.org/abs/math/0703392}{arXiv:math/0703392},
Proposition~6.2 and Corollary~6.3, for the related fixed-base Weil
positivity criterion.
\end{document}
